\section{Introduction}

How a tax system treats the second earner in a couple is a central question in the economics of household labor supply. When spouses are taxed jointly, the lower-earning member's income is stacked on top of the primary earner's, so it faces a high marginal rate from the first hour worked. This discourages secondary-earner participation and can lock in specialisation within the household --- the mechanism that \cite{borella2023} argue holds back female labor supply. Moving to individual taxation lowers the marginal rate on the second income and may therefore raise labor supply, tax revenue, and welfare at the same time.

This assignment studies that reform in a deliberately simple dynamic dual-earner model from lecture~6, stripped of savings, divorce, and uncertainty so that the tax mechanism stands out clearly. We first extend the household objective with a Pareto weight and show why equal weights leave behavior unchanged (Questions~1--2). We then compare joint and individual taxation along three dimensions: life-cycle labor supply (Question~3), the government budget (Questions~4--5), and household welfare under a revenue-neutral reform (Question~6).
